\section{Conclusions}
\label{vi:sec:conclusiones}

Le résultat principal est la construction d’un opérateur positif de masse
sur la fibre d’une particule persistante, avec une spécialisation
scalaire déterminée pour l’électron central. La positivité se déduit de
la transformation par congruence de l’opérateur basal positif, avec un lecteur
autoadjoint et un rapport d’action positif. Sur la trajectoire électronique,
les deux lecteurs produisent \((80,54,6)\) et déterminent la même correction
de retour. Dans la coordinatisation dimensionnelle adoptée, l’évaluation complète donne
\[
 m_e^{\beta}=0.5109989506900008086\ldots\ \mathrm{MeV}/c^2.
\]
Le chiffre réunit le préfacteur électronique et la correction de retour ; sa
confrontation est effectuée avec la valeur de référence et l’incertitude
indiquées dans le tableau~\ref{tab:vi-evaluaciones-veinte}. Le résultat
scalaire spécifique du centre ne remplace ni les opérateurs ni les
conditions d’individualisation des autres collections.

La trajectoire qui soutient cette évaluation n’est pas immobile : sa projection
cellulaire visite exactement les vingt-sept positions définies par
$c-r\equiv0,\pm1\pmod9$. Le registre réversible distingue des histoires qui
partagent un même histogramme de valeurs absolues ; la décomposition spectrale
de $(+++---)^2$ concentre la puissance en $k=2,6,10$ et reproduit exactement
$(\Omega,P_6)=(80,6)$. La
condition opératoire sur la fibre établit en outre qu’un transport TPK
conserve la masse si et seulement s’il conserve la coordonnée complète du caractère et de
l’exposant de retour, condition équivalente à $[M^\beta,V_F]=0$ dans le domaine
stable.

L’espace ambiant régional de \(729\) positions ne reste pas non plus un cardinal
isolé. La bijection par paires de trits
\(\beta:\FF_3^6\to E_9^3\) l’identifie au corps \(S_0\), tandis
que la complétion \(\widehat E^3\) se décompose de manière disjointe en
\(729+243+27+1=1000\) positions. La couronne complète contient \(271\)
positions, l’apex apporte exactement \(1/1000\) de la mesure uniforme et
la compatibilité projective conserve la normalisation en toute dimension.
La preuve et les deux figures du corps exposent séparément la partition
ensembliste et sa réalisation cubique.

La provenance des coefficients est déterminée par l’évolution
ordonnée de \(108\) opérateurs
\(\mathcal U_t=\mathsf{Upd}_t\circ\mathsf{Tra}_t\circ\mathsf{Sel}_t\).
La coordinatisation \(\mathcal R_{12}\) extrait de l’état terminal le registre
incidentiel complet ; \(\mathcal E_{108}^{90,120}\) produit la coordonnée
des canaux ; \(\Pi_H\) extrait la signature ; et \(H_{12}/4\) obtient \(K\).
L’inverse entière restitue exactement les canaux et la somme. Le registre
périodique, la normalisation de \(\alpha\), l’échelle d’action et les
lecteurs de masse agissent sur ces sorties. La même origine structurelle
est donc conservée depuis les opérations du noyau jusqu’à
l’évaluation et sa comparaison métrologique.

\subsection{Identité persistante, atlas et conjugaison}

La persistance est formulée sur des extensions compatibles ; la route
conserve son histoire observable ; la représentation détermine quelles données de
feuillet, d’orientation, de spin, de saveur et de couleur distinguent les états. Cette
identité précède l’opérateur de masse et permet de relier ses lectures
sans réduire la particule à une coordonnée numérique.

L’atlas fournit une classification finie qui peut être parcourue
intégralement. Ses \(81\) cellules, \(56\) collections, \(13\) multisections et
\(324\) conditions orientées sont quatre recensements d’objets différents.
Son organisation identifie l’exception centrale de l’électron et les
fibres non centrales. Dans ces dernières, conserver une multisection ou un
opérateur est une description positive du résultat : elle précise ce qui reste
équivariant et quelle opération supplémentaire produit une lecture spectrale.
L’individualisation physique s’effectue par sa représentation et son
canal, en conservant cette différence.

La conjugaison ne se réduit pas non plus à une unique inversion de signe.
L’inversion cellulaire, l’inversion d’une route, l’involution sur les
mots, la conjugaison de charge et l’holonomie spinorielle conservent
des domaines et des actions propres. La construction de Möbius spécifie la
fibre et le lacet qui engendrent l’inversion d’orientation. Le retour à la
même phase peut laisser une mémoire différente, tandis que la restitution d’une
représentation peut exiger un tour supplémentaire. Ces résultats
décrivent comment la structure persiste lorsque ses observations partielles
coïncident.
La conjugaison antiunitaire précise en outre une distinction dynamique :
la commutation avec l’hamiltonien inverse le sens du groupe
d’évolution, tandis que la préservation de son sous-espace réel à tout instant
exige la relation d’anticommutation correspondante. La compatibilité
avec l’opérateur de masse conserve sa formulation propre.

\subsection{Évaluation massique, composition et occupation}

L’évaluation de masse combine deux niveaux. Le caractère raffiné utilise
la signature entière et les moments des canaux ; la section dimensionnelle
compose action, horloge et coordinatisation des unités. Le préfacteur électronique
reste présent lorsqu’on exprime l’électron par rapport à l’échelle
chronologique. Le rapport des sections d’action agit dans les lecteurs
de retour et permet de comparer les fibres au moyen d’exposants dérivés de
leurs registres. L’égalité des deux lecteurs centraux a reçu
une preuve concrète sur la trajectoire électronique, avec une portée
différenciée des autres profils de l’atlas.
La vitesse de l’horloge et la vitesse constitutive coïncident par
l’identité entre le produit des deux canaux et le déterminant de
l’opérateur de réponse. Lors d’un changement de coordinatisation des unités, leurs coordonnées
se transforment et la grandeur dimensionnelle reste invariante.

La composition des particules a été située avant l’extraction de
la signature. Les registres de liaison conservent les conditions de frontière,
les feuillets et les prédécesseurs sur lesquels agissent les fonctionnelles. Ainsi,
le terme de liaison possède une place définie dans l’opérateur du
composé. Mésons et baryons sont organisés au moyen de leurs espaces de
constituants et de leurs représentations chromatiques ; les résonances
nécessitent en outre la dynamique des canaux ouverts et la définition du
pôle. Leur largeur et leur temps de survie sont des lectures coordonnées
de cette dynamique.

Dans le secteur antipodal neutre, les deux contractions scalaires de
l’extracteur s’annulent exactement :
\(b_{120}^{\nu}=b_{\zeta}^{\nu}=0\). Cette annulation n’élimine pas le
registre orienté, car le sous-espace anti-invariant conserve l’opérateur
complexe \(J^2=-I\) ; elle ne s’étend pas non plus aux routes générales et ne modifie
ni les masses ni le transport PMNS. Dans le prolongement longitudinal, la
comparaison entre dix mots de \(108\) pas et leur concaténation de
\(1080\) pas incorpore le terme de couture explicite. Les nouvelles
frontières ne sont pas des couronnes, de sorte que
\(\Delta C_z=0\) et \(\Delta k_{\Delta_4}=\Delta T_z\), avec \(-2\) comme
valeur atteignable du contrôle local. Ainsi sont conservées séparément la
périodicité de la représentation et la mémoire de l’état TPK complet.

La complétion graduée développe l’occupation sur les modes
complets déjà déclarés. L’antisymétrisation construit les relations
CAR et l’exclusion de Pauli s’en déduit. La symétrisation
construit les CCR sur le domaine des particules en nombre fini et rend visible
le terme de frontière introduit par une troncature bosonique.
La distinction entre égalité d’une masse et coïncidence d’un mode
est exprimée par le déterminant de Gram d’une occupation double.
Spin, feuillet ou couleur peuvent conserver des modes indépendants même lorsqu’une
observable partielle a la même valeur.

Les réalisations orbitales relient l’action aux conditions de
fermeture et à la géométrie de Kepler--Sommerfeld. Leurs hypothèses de métrique,
de dynamique centrale et de représentation ont été maintenues auprès des
équations qui les utilisent. Dans les secteurs de fusion et de tressage sont
explicités l’anneau d’incidence, ses dimensions positives, les
matrices de changement d’association et les caractères d’échange,
ainsi que les identités locales et les conditions de réalisation
que chacun utilise. Cette répartition permet de reconnaître ce qui correspond à
une construction algébrique, à une réalisation dynamique et à une
identification physique.

\Needspace{29\baselineskip}
\subsection{Du spectre à la confrontation métrologique}

La confrontation métrologique complète le parcours au moyen d’observables
typées. Les comparaisons expliquées par secteurs sont accompagnées des
inventaires complets et de leurs sources. Une masse courante, une masse
de pôle, une limite et une largeur conservent leurs définitions, leurs unités
et leurs incertitudes. Chaque évaluation du modèle est reliée à son registre,
à son lecteur et à sa coordinatisation ; chaque observation expérimentale conserve son édition
et sa provenance. Les \(855\) lignes documentaires de masses et de largeurs
constituent le domaine du catalogue de confrontation, distinct du nombre
de prédictions individualisées.

L’apport réunit trois déterminations : l’objet qui persiste,
l’opération qui produit son spectre et la coordinatisation dans laquelle ce spectre est
confronté. La masse résulte de leur composition ; orientation, mémoire et
structure de fibre restent disponibles pour distinguer des états ayant une
même lecture partielle. Le cas électronique réalise cette composition dans
une section scalaire explicite, tandis que l’atlas conserve la diversité
des secteurs opératoires et de leurs réalisations physiques. Cette relation
entre structure et grandeur constitue l’unité du résultat.

% BEGIN SINTESIS_ADITIVA_20260922
\par
La composition supplémentaire ferme la dépendance longitudinale et
gravitationnelle sans modifier le spectre. L’incidence produit
\(L_\alpha=\alpha_{\rm HMT}^{16}(5759/23040)L_\star\) avant de
construire l’horloge ; on en obtient \(R_{108}=L_\alpha\). L’énergie
\(Q_L=\hbar c_{\rm int}/L_\alpha\) fixe conjointement
\(m_{\rm clk}\) et le facteur cotransformé
\(b_e=E_e^{(0)}/Q_L\), de sorte que l’échelle s’annule dans chaque masse déjà
construite. Dans les fibres massives, les deux égalités conservées
\(H\Lambda=\hbar c_{\rm int}I\) et
\(\mathcal R_g\Lambda=L_\alpha^2I\) imposent de manière unique
\(G=c_{\rm int}^3L_\alpha^2/\hbar\). Cette détermination est postérieure à
l’opérateur de masse ; l’expression avec \(t_0\) est son corollaire circulaire et non
son générateur. Le secteur nul conserve son traitement non inversible et
aucune correction décimale dépendant de l’espèce n’est introduite.
% END SINTESIS_ADITIVA_20260922
